% GENERATED FILE. Edit canonical lessons, then run npm run generate:books.
% canonical-source-hash: fnv1a64:8f762debe8ffe6ca
% canonical-lessons: IT-C324-eroe, IT-C324-leggenda, IT-C324-mito, IT-C324-finale, IT-C324-morale

\chapter{Miti ed eroi}
\label{ch:it-miti-ed-eroi}
\hlchaptermodality{324}

\begin{chapteropening}
\textbf{By the end of this chapter:} \emph{I can talk about heroes, legends and myths, how a story ends and its moral.}

Complete the last lesson of chapter 324: I can talk about heroes, legends and myths, how a story ends and its moral.
\end{chapteropening}

\section[l'eroe]{l'eroe — the hero}

\label{lesson:IT-C324-eroe}



\begin{quote}
Before the new one: say the Italian for to narrate, then the Italian for the narrator.
\end{quote}

\subsection*{You'll want to know: l'eroe}
\textbf{l'eroe} — \textquotedblleft{}the hero\textquotedblright{}.

\begin{quote}
L'eroe della storia è un ragazzo povero.
\end{quote}

\begin{quote}
\emph{The hero of the story is a poor boy.}
\end{quote}

A woman is \emph{l'eroina}.

\begin{cousinweb}
\emph{eroe} is a learned borrowing from Latin \emph{hērōs}, from Greek \emph{h\'{\={e}}rōs} \textquotedblleft{}a hero, a demigod\textquotedblright{} (Lewis \& Short).
\end{cousinweb}

\begin{grammarlens}[title={What you've built}]
One more word for telling, explaining and describing at B1.
\end{grammarlens}

\subsection*{Your turn}
\begin{itemize}
\raggedright
  \item \emph{Say it:} \emph{l'eroe}
  \item \emph{Say it:} \emph{l'eroe}, once more
  \item \emph{Say it:} a sentence of your own with \emph{l'eroe}
  \item \emph{Recall:} say the Italian for to narrate, then the Italian for the narrator, then say \emph{l'eroe} again
\end{itemize}

\begin{tcolorbox}[breakable,colback=teal!4,colframe=teal!35!black,title={Before you move on}]
What does l'eroe mean? (\textquotedblleft{}The hero\textquotedblright{}.) Say it once more.
\end{tcolorbox}
\section[la leggenda]{la leggenda — the legend}

\label{lesson:IT-C324-leggenda}



\begin{quote}
Before the new one: say the Italian for the narrator, then the Italian for the hero.
\end{quote}

\subsection*{You'll want to know: la leggenda}
\textbf{la leggenda} — \textquotedblleft{}the legend\textquotedblright{}.

\begin{quote}
Secondo la leggenda, la città è nata in una notte.
\end{quote}

\begin{quote}
\emph{According to the legend, the city was born in a single night.}
\end{quote}

\emph{Una leggenda vivente} — a living legend.

\begin{cousinweb}
\emph{leggenda} comes from Medieval Latin \emph{legenda} \textquotedblleft{}things to be read\textquotedblright{}, from \emph{legere} \textquotedblleft{}to read\textquotedblright{} (Lewis \& Short): first, the lives of saints read aloud.
\end{cousinweb}

\begin{grammarlens}[title={What you've built}]
One more word for telling, explaining and describing at B1.
\end{grammarlens}

\subsection*{Your turn}
\begin{itemize}
\raggedright
  \item \emph{Say it:} \emph{la leggenda}
  \item \emph{Say it:} \emph{la leggenda}, once more
  \item \emph{Say it:} a sentence of your own with \emph{la leggenda}
  \item \emph{Recall:} say the Italian for the narrator, then the Italian for the hero, then say \emph{la leggenda} again
\end{itemize}

\begin{tcolorbox}[breakable,colback=teal!4,colframe=teal!35!black,title={Before you move on}]
What does la leggenda mean? (\textquotedblleft{}The legend\textquotedblright{}.) Say it once more.
\end{tcolorbox}
\section[il mito]{il mito — the myth}

\label{lesson:IT-C324-mito}



\begin{quote}
Before the new one: say the Italian for the hero, then the Italian for the legend.
\end{quote}

\subsection*{You'll want to know: il mito}
\textbf{il mito} — \textquotedblleft{}the myth\textquotedblright{}.

\begin{quote}
Questo è un mito molto antico.
\end{quote}

\begin{quote}
\emph{This is a very old myth.}
\end{quote}

Also a hero of today: \emph{quel cantante è un mito!}

\begin{cousinweb}
\emph{mito} comes from Greek \emph{mŷthos} \textquotedblleft{}a word, a tale\textquotedblright{} (EtymDB), through Late Latin \emph{mythos}.
\end{cousinweb}

\begin{grammarlens}[title={What you've built}]
One more word for telling, explaining and describing at B1.
\end{grammarlens}

\subsection*{Your turn}
\begin{itemize}
\raggedright
  \item \emph{Say it:} \emph{il mito}
  \item \emph{Say it:} \emph{il mito}, once more
  \item \emph{Say it:} a sentence of your own with \emph{il mito}
  \item \emph{Recall:} say the Italian for the hero, then the Italian for the legend, then say \emph{il mito} again
\end{itemize}

\begin{tcolorbox}[breakable,colback=teal!4,colframe=teal!35!black,title={Before you move on}]
What does il mito mean? (\textquotedblleft{}The myth\textquotedblright{}.) Say it once more.
\end{tcolorbox}
\section[il finale]{il finale — the ending (of a story or film)}

\label{lesson:IT-C324-finale}



\begin{quote}
Before the new one: say the Italian for the legend, then the Italian for the myth.
\end{quote}

\subsection*{You'll want to know: il finale}
\textbf{il finale} — \textquotedblleft{}the ending (of a story or film)\textquotedblright{}.

\begin{quote}
Il finale del film mi ha fatto piangere.
\end{quote}

\begin{quote}
\emph{The ending of the film made me cry.}
\end{quote}

\emph{Il lieto fine} is a happy ending.

\begin{cousinweb}
\emph{finale} is inherited from Late Latin \emph{fīnālis} (EtymDB), built on \emph{fīnis} \textquotedblleft{}end\textquotedblright{}.
\end{cousinweb}

\begin{grammarlens}[title={What you've built}]
One more word for telling, explaining and describing at B1.
\end{grammarlens}

\subsection*{Your turn}
\begin{itemize}
\raggedright
  \item \emph{Say it:} \emph{il finale}
  \item \emph{Say it:} \emph{il finale}, once more
  \item \emph{Say it:} a sentence of your own with \emph{il finale}
  \item \emph{Recall:} say the Italian for the legend, then the Italian for the myth, then say \emph{il finale} again
\end{itemize}

\begin{tcolorbox}[breakable,colback=teal!4,colframe=teal!35!black,title={Before you move on}]
What does il finale mean? (\textquotedblleft{}The ending (of a story or film)\textquotedblright{}.) Say it once more.
\end{tcolorbox}
\section[la morale]{la morale — the moral (of a story)}

\label{lesson:IT-C324-morale}



\begin{quote}
Before the new one: say the Italian for the myth, then the Italian for the ending.
\end{quote}

\subsection*{You'll want to know: la morale}
\textbf{la morale} — \textquotedblleft{}the moral (of a story)\textquotedblright{}.

\begin{quote}
La morale della storia è semplice.
\end{quote}

\begin{quote}
\emph{The moral of the story is simple.}
\end{quote}

Also morals in general. Masculine \emph{il morale} is your spirits.

\begin{cousinweb}
\emph{morale} is a learned borrowing from Latin \emph{mōrālis} \textquotedblleft{}of manners\textquotedblright{} (EtymDB), built on \emph{mōs, mōris} \textquotedblleft{}custom\textquotedblright{}.
\end{cousinweb}

\begin{grammarlens}[title={What you've built}]
One more word for telling, explaining and describing at B1.
\end{grammarlens}

\subsection*{Your turn}
\begin{itemize}
\raggedright
  \item \emph{Say it:} \emph{la morale}
  \item \emph{Say it:} \emph{la morale}, once more
  \item \emph{Say it:} a sentence of your own with \emph{la morale}
  \item \emph{Recall:} say the Italian for the myth, then the Italian for the ending, then say \emph{la morale} again
\end{itemize}

\begin{tcolorbox}[breakable,colback=teal!4,colframe=teal!35!black,title={Before you move on}]
What does la morale mean? (\textquotedblleft{}The moral (of a story)\textquotedblright{}.) Say it once more.
\end{tcolorbox}
